\section{Revelation, Tradition, Apparition, and Legend}

The word ``revelation'' can name radically different things. Catholic doctrine therefore requires a hierarchy of claims, especially in a devotion so richly surrounded by saintly traditions and apparitions.

\subsection{The definitive Revelation in Christ}

God's public Revelation reaches its fullness in the incarnate Son. No new public revelation is to be expected before Christ's glorious manifestation, and no later message can correct or complete the deposit entrusted to the Apostles (Second Vatican Council, \emph{Dei Verbum} 2--4; CCC 65--66). The Rosary's primary revealed content is this public economy: the scriptural events, the apostolic confession of Christ, and doctrines authoritatively transmitted by Scripture and Tradition.

The prayer form itself is not a new article of faith. Its scriptural words, ecclesial petitions, selection of mysteries, and method developed in the Church. Its authority is real but differentiated: Gospel and dogma bind as divine and Catholic faith requires; papal commendation and ecclesial norms guide devotion according to their proper weight; optional customary forms remain optional.

\subsection{Private revelations}

So-called private revelations may occur after the apostolic age. Their role is not to improve Christ's definitive Revelation but to help the faithful live it more fully in a particular historical situation. Ecclesial judgments must be read according to the act that issued them. Earlier recognitions could judge an alleged apparition worthy of belief or even affirm supernatural origin. Under the Dicastery for the Doctrine of the Faith's 2024 norms, the ordinary positive conclusion is a \emph{nihil obstat}: it permits, and may encourage, positive pastoral reception without certifying that the alleged phenomenon is supernatural. No private revelation enters the deposit, becomes an object of faith, or requires theological faith from every Catholic (CCC 67; DDF, \emph{Norms for Proceeding in the Discernment of Alleged Supernatural Phenomena} 11--12, 17, 23).

The Rosary appears prominently in ecclesially received apparitions, above all Lourdes and Fatima. At Fatima the reported requests for daily Rosary prayer, penance, conversion, and peace intensify themes already belonging to the Gospel and the Church's piety. The Shrine's official narrative locates the repeated Rosary requests in the reported encounters from May through October 1917 and the commonly used post-decade invocation in the July account. These reports do not originate the Hail Mary, the decades, or the traditional mysteries. The prayer commonly called the Fatima prayer is a widely used optional invocation after a decade; it is not part of the Rosary's universally stable core and is not required for a five-decade recitation.

Even an approved apparition is mediated through human perception, memory, language, and later reporting. The Congregation for the Doctrine of the Faith's theological commentary on Fatima distinguishes the prophetic symbolism of private revelation from a photographic preview of fixed future events. Interpretation must remain Christological, ecclesial, sober, and open to the prudence of legitimate authority.

\subsection{The Dominic claim classified}

Three propositions should not be merged:

\begin{enumerate}
  \item The Dominican family had a major historical role in preaching, organizing, and defending the Marian Psalter and Rosary.
  \item Catholic devotional and papal tradition has associated St.~Dominic personally with Marian inspiration and the Rosary.
  \item The complete modern Rosary was dictated in its final form to Dominic in a historically documented apparition.
\end{enumerate}

The first is historically well grounded; the second is a real strand of received devotional tradition; the third is not established by contemporary documentation and does not fit the observable development of the prayer's components. Respect for tradition does not require pretending that these propositions have equal evidentiary status. Indeed, MC 38--39 explicitly commands Marian devotion to avoid vain credulity, falsehood, and the search for extraordinary novelty.

\subsection{Doctrine, typology, meditation, and imagination}

Four further levels recur within the mystery studies:

\begin{threestable}{Level}{Example}{Proper claim}
Defined or authoritative doctrine & Incarnation, Christ's two wills, real Resurrection, Eucharistic presence, Assumption. & Receive according to the authority of the source; do not reduce dogma to symbol.\\
Biblical and traditional typology & Mary as new Eve or Ark; the Cross from the new Adam's side; thorns and the curse; Pentecost and Babel. & A real theological correspondence in the one economy, but not always the literal sense or a narrated historical detail.\\
Devotional meditation & Imagining the poverty of Bethlehem or standing with Mary at Calvary. & Can engage affection and moral response if disciplined by Scripture and doctrine. It is not recovered eyewitness data.\\
Unverified or legendary detail & Exact numbers of wounds, alleged objects, undocumented speeches, or a late narrative presented as contemporary evidence. & May not be asserted as fact without adequate evidence; must never become a condition of faith.\\
\end{threestable}

The imagination is not an enemy of contemplation. The Incarnation makes places, gestures, voices, bodies, and relationships spiritually significant. But imagination must be hospitable rather than sovereign: it serves the revealed event, accepts correction, and does not fill silence with sensational certainty.

\subsection{A rule of ecclesial discernment}

A Rosary-related message or practice should be assessed by whether it confesses the Catholic faith, leads to Christ and the sacraments, produces humility and charity, respects ecclesial authority and Christian freedom, avoids commercial or manipulative exploitation, and remains truthful about its evidence. Claims that generate contempt for the Church, date the end, promise immunity through an object, make salvation depend on secret knowledge, or turn Mary against her Son contradict the prayer's own structure.

The safest Marian maxim remains Cana's: do what Christ tells you. The Church does not honor Mary by granting an alleged message authority she herself never possesses independently of God.
